% ---------------------------------------------------------------------
%  05 Experimental Setup
%  This chapter describes *how* things are measured. Numbers go in 06.
% ---------------------------------------------------------------------
\section{Experimental Setup}
\label{sec:experiments}

\subsection{Datasets}
\label{subsec:datasets}
\Todo{One paragraph per dataset: source, number of clips/frames, annotation
procedure, train/val/test split policy, and whether the split is by clip or by
frame (say it explicitly --- frame-level splits leak).}

\begin{table}[t]
  \centering
  \caption{\Todo{Dataset summary: clips, frames, cameras, annotations.}}
  \label{tab:datasets}
  % \input{assets/tables/datasets.tex}
  \begin{tabular}{lrrr}
    \toprule
    Dataset & Clips & Frames & Views \\
    \midrule
    \Todo{fill} & -- & -- & -- \\
    \bottomrule
  \end{tabular}
\end{table}

\subsection{Metrics}
\label{subsec:metrics}
\Todo{Define each metric formally, with units. At minimum:
2D ball localisation error and detection rate;
court keypoint PCK and homography reprojection error;
3D ball trajectory error (m);
\PLCS{} root position error (m) and rotation error (deg);
end-to-end court-frame joint error.
State which metrics are computed only on visible/valid frames.}

\subsection{Baselines}
\label{subsec:baselines}
\Todo{What we compare against and why each comparison is fair. Include the
geometric baseline (direct triangulation from the homographies) --- it is the
honest lower bound the learned lifters must beat.}

\subsection{Ablations}
\label{subsec:ablation_protocol}
\Todo{List the ablation axes and the single question each one answers. One
knob per row.}

\subsection{Evaluation protocol}
\label{subsec:protocol}
\Todo{Seeds and number of runs; what is reported (mean $\pm$ std vs. best);
checkpoint selection criterion. Reproduction bundles are described in
\cref{sec:appendix_repro}.}
